\documentclass[11pt,a4paper]{article}
\usepackage[T1]{fontenc}
\usepackage[utf8]{inputenc}
\usepackage[main=italian,provide=*]{babel}
\usepackage{amsmath,amssymb}
\usepackage{geometry}
\geometry{margin=2.3cm,top=2.5cm,bottom=2.7cm}
\usepackage{booktabs}
\usepackage[table]{xcolor}
\usepackage{tcolorbox}
\tcbuselibrary{skins,breakable}
\usepackage{titlesec}
\usepackage{enumitem}
\usepackage{seqsplit}
\usepackage{needspace}
\usepackage{fancyhdr}
\usepackage{hyperref}
\hypersetup{colorlinks=true,linkcolor=Sepia,citecolor=Sepia,urlcolor=Sepia}

\definecolor{inkblue}{HTML}{1B3A5C}
\definecolor{lightblue}{HTML}{EAF1F8}
\definecolor{accent}{HTML}{B0562A}
\definecolor{softgray}{HTML}{6B6B6B}
\definecolor{okgreen}{HTML}{2E6B3E}
\definecolor{okgreenbg}{HTML}{EAF5EC}
\definecolor{warnamber}{HTML}{9C6B12}
\definecolor{warnbg}{HTML}{FBF1E0}
\definecolor{advisorpurple}{HTML}{5B3A7A}
\definecolor{advisorbg}{HTML}{F1ECF6}
\definecolor{notebar}{HTML}{9C6B12}

\titleformat{\section}
  {\normalfont\Large\bfseries\color{inkblue}}
  {\thesection}{0.6em}{}[{\color{inkblue!40}\titlerule}]
\titlespacing{\section}{0pt}{0.8em}{0.4em}
\titleformat{\subsection}
  {\normalfont\large\bfseries\color{accent}}
  {}{0em}{}
\titlespacing{\subsection}{0pt}{0.5em}{0.2em}

\pagestyle{fancy}
\fancyhf{}
\renewcommand{\headrulewidth}{0.4pt}
\fancyhead[R]{\small\color{softgray}\thepage}
\fancyfoot[C]{\small\color{softgray}Tesi triennale, Samuele Schiavi --- Relatore: Prof.\ Alessandro Chiesa}

\newtcolorbox{keyeq}[1][]{
  colback=lightblue, colframe=inkblue!70, boxrule=0.6pt,
  arc=2pt, left=8pt, right=8pt, top=4pt, bottom=4pt, #1
}
\newtcolorbox{panelbox}[1]{
  colback=white, colframe=softgray!50, boxrule=0.5pt,
  arc=2pt, left=7pt, right=7pt, top=3pt, bottom=3pt,
  fonttitle=\bfseries\color{inkblue}, title={#1}
}
\newtcolorbox{okbox}[1]{
  colback=okgreenbg, colframe=okgreen!60, boxrule=0.6pt,
  arc=2pt, left=7pt, right=7pt, top=3pt, bottom=3pt,
  fonttitle=\bfseries\color{okgreen}, title={#1}
}
\newtcolorbox{warnbox}[1]{
  colback=warnbg, colframe=warnamber!70, boxrule=0.6pt,
  arc=2pt, left=7pt, right=7pt, top=3pt, bottom=3pt,
  fonttitle=\bfseries\color{warnamber}, title={#1}
}
\newtcolorbox{advisorbox}[1]{
  colback=advisorbg, colframe=advisorpurple!60, boxrule=0.6pt,
  arc=2pt, left=7pt, right=7pt, top=4pt, bottom=4pt,
  fonttitle=\bfseries\color{advisorpurple}, title={#1}
}
\newtcolorbox{editnote}{
  colback=white, colframe=white, boxrule=0pt,
  borderline west={2.2pt}{0pt}{notebar},
  arc=0pt, left=10pt, right=4pt, top=2pt, bottom=2pt,
  fontupper=\itshape\small\color{softgray!90!black}
}
\fancyhead[L]{\small\color{softgray}Domande per il relatore --- quantum simulation}

\title{\vspace{-1.5em}\color{inkblue}Domande per il relatore\\[0.1em]\Large\normalfont\color{softgray}Quantum simulation (Trotter) sul dimero}
\author{Note di lavoro --- Tesi triennale, Samuele Schiavi\\ Relatore: Prof.\ Alessandro Chiesa}
\date{}

\begin{document}
\sloppy
\maketitle
\thispagestyle{fancy}
\vspace{-2em}

Domande aperte sul nuovo task (simulare $e^{-iHt}$ sull'Hamiltoniana
del dimero, non su Ising), da porre a Prof.\ Chiesa al prossimo
incontro.

\section{Serve un qubit di supporto per le correlazioni nel tempo?}

Nel notebook di esempio che ho usato come riferimento (quello sul
modello di Ising), oltre ai qubit che rappresentano i due spin fisici
ne viene usato un \textbf{terzo, aggiuntivo}, che non rappresenta
nessuno spin reale del sistema: serve solo come "strumento di misura
indiretto" per calcolare quantità che altrimenti non sarebbero
accessibili direttamente, tipo la correlazione tra uno spin al tempo
$t$ e l'altro al tempo $0$, $\langle s_{x1}(t)s_{x2}\rangle$.

\begin{advisorbox}{Domanda}
Per il task che mi ha chiesto (evoluzione temporale dell'Hamiltoniana
del dimero, non del modello di Ising), mi serve anche a me questo
qubit aggiuntivo, oppure per ora mi basta guardare l'evoluzione di
osservabili "semplici" (es.\ la magnetizzazione $M_z(t)$), senza
calcolare correlazioni a due tempi?
\end{advisorbox}

\section{Includo subito il termine DM o parto dal caso più semplice?}

Con il termine DM acceso ($D\neq0$) i due pezzi dell'Hamiltoniana (il
campo $b$ e lo scambio $J$) non commutano, quindi serve una
decomposizione approssimata (Trotter), con un errore che dipende dal
numero di passi usati.

Se invece $D=0$, i due pezzi commutano e si può scrivere l'evoluzione
in modo \textbf{esatto}, senza approssimazioni.

\begin{advisorbox}{Domanda}
Conviene partire dal caso $D=0$ (esatto) per validare il metodo, e poi
passare a $D\neq0$ (Trotter, con l'errore da studiare)? O vuole vedere
subito il caso completo con $D\neq0$?
\end{advisorbox}

\section{Dove va collocato questo blocco nella tesi?}

\begin{advisorbox}{Domanda}
Questa simulazione temporale è una sotto-sezione della Parte 1
(sistema chiuso), oppure la immagina come una sezione a sé, magari come
ponte verso la Parte 2 (sistema aperto/rumore)?
\end{advisorbox}

\section{Range di parametri e condizioni iniziali}

\begin{advisorbox}{Domanda}
C'è un valore di riferimento per $b/J$ e $D$ da usare (es.\ gli stessi
usati nel resto della tesi per il dimero), e uno stato iniziale
specifico da cui far partire l'evoluzione (es.\
$|\!\uparrow\uparrow\rangle$, oppure lo stato fondamentale trovato con
il VQE)?
\end{advisorbox}

\end{document}
